\documentclass[10pt,a4paper]{amsart}
\usepackage{amsmath,amssymb,amsthm,graphicx}
\usepackage{verbatim}
\newenvironment{reference}{\comment}{\endcomment}
%\newenvironment{reference}{}{}
\newenvironment{slogan}{\comment}{\endcomment}
\newenvironment{history}{\comment}{\endcomment}
% For commutative diagrams we use Xy-pic
\usepackage[all]{xy}

% We use 2cell for 2-commutative diagrams.
\xyoption{2cell}
% Theorem environments.
%
\theoremstyle{plain}
\newtheorem{theorem}[subsection]{Theorem}
\newtheorem{proposition}[subsection]{Proposition}
\newtheorem{lemma}[subsection]{Lemma}

\theoremstyle{definition}
\newtheorem{definition}[subsection]{Definition}
\newtheorem{example}[subsection]{Example}
\newtheorem{exercise}[subsection]{Exercise}
\newtheorem{situation}[subsection]{Situation}

\theoremstyle{remark}
\newtheorem{remark}[subsection]{Remark}
\newtheorem{remarks}[subsection]{Remarks}

\numberwithin{equation}{subsection}

% Macros
%
\def\lim{\mathop{\mathrm{lim}}\nolimits}
\def\colim{\mathop{\mathrm{colim}}\nolimits}
\def\Spec{\mathop{\mathrm{Spec}}}
\def\Hom{\mathop{\mathrm{Hom}}\nolimits}
\def\Ext{\mathop{\mathrm{Ext}}\nolimits}
\def\SheafHom{\mathop{\mathcal{H}\!\mathit{om}}\nolimits}
\def\SheafExt{\mathop{\mathcal{E}\!\mathit{xt}}\nolimits}
\def\Sch{\mathit{Sch}}
\def\Mor{\mathop{\mathrm{Mor}}\nolimits}
\def\Ob{\mathop{\mathrm{Ob}}\nolimits}
\def\Sh{\mathop{\mathit{Sh}}\nolimits}
\def\NL{\mathop{N\!L}\nolimits}
\def\CH{\mathop{\mathrm{CH}}\nolimits}
\def\proetale{{pro\text{-}\acute{e}tale}}
\def\etale{{\acute{e}tale}}
\def\QCoh{\mathit{QCoh}}
\def\Ker{\mathop{\mathrm{Ker}}}
\def\Im{\mathop{\mathrm{Im}}}
\def\Coker{\mathop{\mathrm{Coker}}}
\def\Coim{\mathop{\mathrm{Coim}}}

% Boxtimes
%
\DeclareMathSymbol{\boxtimes}{\mathbin}{AMSa}{"02}

%
% Macros for moduli stacks/spaces
%
\def\QCohstack{\mathcal{QC}\!\mathit{oh}}
\def\Cohstack{\mathcal{C}\!\mathit{oh}}
\def\Spacesstack{\mathcal{S}\!\mathit{paces}}
\def\Quotfunctor{\mathrm{Quot}}
\def\Hilbfunctor{\mathrm{Hilb}}
\def\Curvesstack{\mathcal{C}\!\mathit{urves}}
\def\Polarizedstack{\mathcal{P}\!\mathit{olarized}}
\def\Complexesstack{\mathcal{C}\!\mathit{omplexes}}
% \Pic is the operator that assigns to X its picard group, usage \Pic(X)
% \Picardstack_{X/B} denotes the Picard stack of X over B
% \Picardfunctor_{X/B} denotes the Picard functor of X over B
\def\Pic{\mathop{\mathrm{Pic}}\nolimits}
\def\Picardstack{\mathcal{P}\!\mathit{ic}}
\def\Picardfunctor{\mathrm{Pic}}
\def\Deformationcategory{\mathcal{D}\!\mathit{ef}}
\usepackage{cleveref}
\crefname{originallemma}{original lemma}{original lemmas}
\Crefname{originallemma}{Original lemma}{Original lemmas}
\makeatletter
\newlabel{lemma-descent-smooth}{{lemma-descent-smooth}{source}{}{sourcepointer-1}{}}
\newlabel{lemma-descent-smooth@cref}{{[originallemma][0][]lemma-descent-smooth}{[1][0][]source}{}{}{}}
\makeatother
\crefname{originallemma}{original lemma}{original lemmas}
\Crefname{originallemma}{Original lemma}{Original lemmas}
\makeatletter
\newlabel{lemma-gerbe-fppf}{{lemma-gerbe-fppf}{source}{}{sourcepointer-2}{}}
\newlabel{lemma-gerbe-fppf@cref}{{[originallemma][0][]lemma-gerbe-fppf}{[1][0][]source}{}{}{}}
\makeatother
\crefname{originallemma}{original lemma}{original lemmas}
\Crefname{originallemma}{Original lemma}{Original lemmas}
\makeatletter
\newlabel{lemma-local-structure-gerbe}{{lemma-local-structure-gerbe}{source}{}{sourcepointer-3}{}}
\newlabel{lemma-local-structure-gerbe@cref}{{[originallemma][0][]lemma-local-structure-gerbe}{[1][0][]source}{}{}{}}
\makeatother
\crefname{originallemma}{original lemma}{original lemmas}
\Crefname{originallemma}{Original lemma}{Original lemmas}
\makeatletter
\newlabel{lemma-smooth-quotient-stack}{{lemma-smooth-quotient-stack}{source}{}{sourcepointer-4}{}}
\newlabel{lemma-smooth-quotient-stack@cref}{{[originallemma][0][]lemma-smooth-quotient-stack}{[1][0][]source}{}{}{}}
\makeatother
\title{A gerbe whose total space is smooth over the base has a smooth structure morphism}
\author{The Stacks Project Authors}
\date{stacks-project:stacks-morphisms.tex:local-snapshot-20261012; GNU FDL 1.2 or later}
\begin{document}
\maketitle
\noindent\textit{Editorial note.} All delivered mathematical text is the author's own, copied from the pinned Stacks Project chapter source sources/stacks-algebraic-stacks-morphisms/pinned-original-source.tex; no statement or proof step was rewritten, reordered or removed.\par\medskip
\noindent\textit{Editorial note.} Preamble provenance: the Stacks Project keeps one single shared preamble for every chapter; this card takes it directly from the staged sources/stacks-algebraic-stacks-morphisms/preamble.tex as three verbatim spans (lines 8-12, 14-18 and 154-217), so no derived file was written and nothing was added or redefined. The wrapper class and the base packages are supplied by the assembler.\par\medskip
\noindent\textit{Editorial note.} Bibliography provenance: neither the counted statement nor its proof contains a citation (this lane rescans the whole span before writing the card), so the card carries no bibliography and no entry was reconstructed.\par\medskip
\noindent\textit{Editorial note.} Selection provenance: this card was opened under the file-based control contract (CONTROL.md), with the mechanical pre-card gate run on its statement and proof interval before the card was written.\par\medskip
\noindent\textit{Allowed standard prerequisites (editorial).}\begin{itemize}
\item The statement that the structure morphism of a gerbe is flat and locally of finite presentation. Invoked at the start of the counted proof (source line 7734). Supplies one half of the smoothness assertion.
\item The descent statement for smoothness along a flat finitely presented cover. Invoked inside the counted proof (source lines 7735-7737). Reduces the assertion to a flat finitely presented base change.
\item The local structure of a gerbe as a quotient stack over an algebraic space. Invoked inside the counted proof (source lines 7738-7741). Presents the base-changed gerbe as a quotient by a flat finitely presented group.
\item The statement that a quotient of a smooth space by a flat finitely presented group is smooth over the base. Invoked at the end of the counted proof (source line 7744). Gives smoothness in the quotient case.
\end{itemize}\medskip
\ifdefined\hypertarget\hypertarget{sourcepointer-1}{}\fi
\noindent\textit{Source reference.} Original lemma lemma-descent-smooth: sources/stacks-algebraic-stacks-morphisms/pinned-original-source.tex, source lines 7627--7627. Navigation pointer only: the counted text refers to the Stacks item lemma-descent-smooth, which is an existing labeled result of this same chapter and is not part of this card. The pointer carries the target's own label as its number and its definition line as its span; its content is not reproduced here.\par\medskip
\ifdefined\hypertarget\hypertarget{sourcepointer-2}{}\fi
\noindent\textit{Source reference.} Original lemma lemma-gerbe-fppf: sources/stacks-algebraic-stacks-morphisms/pinned-original-source.tex, source lines 6287--6287. Navigation pointer only: the counted text refers to the Stacks item lemma-gerbe-fppf, which is an existing labeled result of this same chapter and is not part of this card. The pointer carries the target's own label as its number and its definition line as its span; its content is not reproduced here.\par\medskip
\ifdefined\hypertarget\hypertarget{sourcepointer-3}{}\fi
\noindent\textit{Source reference.} Original lemma lemma-local-structure-gerbe: sources/stacks-algebraic-stacks-morphisms/pinned-original-source.tex, source lines 6248--6248. Navigation pointer only: the counted text refers to the Stacks item lemma-local-structure-gerbe, which is an existing labeled result of this same chapter and is not part of this card. The pointer carries the target's own label as its number and its definition line as its span; its content is not reproduced here.\par\medskip
\ifdefined\hypertarget\hypertarget{sourcepointer-4}{}\fi
\noindent\textit{Source reference.} Original lemma lemma-smooth-quotient-stack: sources/stacks-algebraic-stacks-morphisms/pinned-original-source.tex, source lines 7698--7698. Navigation pointer only: the counted text refers to the Stacks item lemma-smooth-quotient-stack, which is an existing labeled result of this same chapter and is not part of this card. The pointer carries the target's own label as its number and its definition line as its span; its content is not reproduced here.\par\medskip
\begin{lemma}
\label{lemma-gerbe-smooth}
Let $\pi : \mathcal{X} \to \mathcal{Y}$ be a morphism of algebraic stacks.
If $\mathcal{X}$ is a gerbe over $\mathcal{Y}$, then $\pi$ is surjective
and smooth.
\end{lemma}

\begin{proof}
We have seen surjectivity in Lemma \ref{lemma-gerbe-fppf}.
By Lemma \ref{lemma-descent-smooth}
it suffices to prove to the lemma after replacing $\pi$ by a base change
with a surjective, flat, locally finitely presented morphism
$\mathcal{Y}' \to \mathcal{Y}$. By
Lemma \ref{lemma-local-structure-gerbe}
we may assume $\mathcal{Y} = U$ is an algebraic space and
$\mathcal{X} = [U/G]$ over $U$ with $G \to U$ flat and
locally of finite presentation.
Then we win by Lemma \ref{lemma-smooth-quotient-stack}.
\end{proof}

\end{document}
